\renewcommand\bibname{参考文献}
\begin{thebibliography}{99}
	
	\bibitem{aho2006}
	Alfred V.~Aho, Monica S.~Lam, Ravi~Sethi, and Jeffrey D.~Ullman.
	\newblock {\em Compilers: Principles, Techniques, and Tools},
	2nd edition.
	\newblock Addison--Wesley, 2006.
	
	\bibitem{yiu2014}
	Joseph~Yiu.
	\newblock {\em The Definitive Guide to ARM Cortex-M3 and Cortex-M4 Processors},
	3rd edition.
	\newblock Newnes, 2014.
	
	\bibitem{douglass2011}
	Bruce Powel~Douglass.
	\newblock {\em Design Patterns for Embedded Systems in C}.
	\newblock Elsevier, 2011.
	
	\bibitem{chen2020}
	陈海波，夏虞斌等.
	\newblock {\em 嵌入式系统设计：CPS与物联网物理世界}.
	\newblock 高等教育出版社，2020.
	
	\bibitem{niaoge}
	鸟哥.
	\newblock {\em 鸟哥的Linux私房菜（基础学习篇）}.
	\newblock 人民邮电出版社.
	
	\bibitem{blum2015}
	Richard~Blum and Christine~Bresnahan.
	\newblock {\em Linux命令行与Shell脚本编程大全}（第3版）.
	\newblock 人民邮电出版社，2015.
	
	\bibitem{liuchan}
	刘遄.
	\newblock {\em Linux就该这么学}.
	\newblock 人民邮电出版社.
	
	\bibitem{stevens2013}
	W.~Richard~Stevens and Stephen A.~Rago.
	\newblock {\em UNIX环境高级编程}（第3版）.
	\newblock 人民邮电出版社，2013.
	
	\bibitem{zhuwenwei}
	朱文伟，李建英.
	\newblock {\em Linux C与C++一线开发实践}.
	\newblock 清华大学出版社.
	
	\bibitem{songjingbin}
	宋敬彬.
	\newblock {\em Linux网络编程}（第2版）.
	\newblock 清华大学出版社.
	
	\bibitem{weidongshan}
	韦东山.
	\newblock {\em 嵌入式Linux应用开发完全手册}.
	\newblock 人民邮电出版社.
	
	\bibitem{corbet2005}
	Jonathan~Corbet, Alessandro~Rubini, and Greg~Kroah-Hartman.
	\newblock {\em Linux设备驱动程序}（第3版）.
	\newblock O'Reilly, 2005.
	
	\bibitem{songbaohua}
	宋宝华.
	\newblock {\em Linux设备驱动开发详解}（第2版）.
	\newblock 人民邮电出版社.
	
	\bibitem{madieu}
	John~Madieu.
	\newblock {\em Linux设备驱动开发}（第2版）.
	\newblock 人民邮电出版社.
	
	\bibitem{bovet2005}
	Daniel P.~Bovet and Marco~Cesati.
	\newblock {\em 深入理解Linux内核}（第3版）.
	\newblock 中国电力出版社，2005.
	
	\bibitem{love2010}
	Robert~Love.
	\newblock {\em Linux内核设计与实现}（原书第3版，典藏版）.
	\newblock 机械工业出版社，2010.
	
	\bibitem{jiangyahua}
	姜亚华.
	\newblock {\em 图解Linux内核：基于6.x}.
	\newblock 机械工业出版社，2024.
	
	\bibitem{zhang2026}
	最爱小兔糖糖.
	\newblock {\em 基于STM32和边缘视觉的智能小车设计与实现}.
	\newblock 合肥：安徽新华学院，2026.
	
\end{thebibliography}